\documentclass[11pt,a4paper]{article}
\usepackage[utf8]{inputenc}
\usepackage[T1]{fontenc}
\usepackage[english]{babel}
\usepackage{graphicx}
\usepackage{hyperref}
\usepackage{booktabs}
\usepackage{geometry}
\geometry{margin=2.5cm}
\usepackage{xcolor}
\usepackage{verbatim}
\usepackage{parskip}
\setlength{\parindent}{0pt}
\title{ALC Ecosystem Graph Analysis Pipeline}
\author{ALC}
\date{\today}
\begin{document}
\maketitle
\tableofcontents
\newpage
\section{Introduction}

\subsection{An Introduction to K-Tool}

K-Tool is a digital platform developed by the Agirre Lehendakaria Center
(ALC) to support systemic innovation. It enables teams to collect,
organize, and analyze information about complex social ecosystems
throughout ALC's methodology.

The platform supports four main phases:

\begin{itemize}
    \item \textbf{Mapping} -- Identify the ecosystem and its actors.
    \item \textbf{Listening} -- Collect qualitative information and identify shared
\end{itemize}
  perceptions, challenges, and opportunities.
\begin{itemize}
    \item \textbf{Sensemaking} -- Interpret the collected information to understand
\end{itemize}
  systemic dynamics.
\begin{itemize}
    \item \textbf{Co-creation} -- Design initiatives that respond to the identified
\end{itemize}
  challenges.

K-Tool stores a wide range of interconnected information, including
actors, organizations, projects, interviews, quotes, perceptions,
challenges, opportunities, values, documents, and the relationships
between them. While this creates a valuable knowledge base, several
limitations prevent the platform from functioning as a true social
digital twin.

\subsection{Current Limitations}

Although K-Tool contains AI-assisted features and a rich collection of
interconnected information, its underlying architecture remains
fundamentally record-based. Data are distributed across multiple
relational tables, meaning relationships exist only where they have been
explicitly entered by users. Consequently, the platform cannot model the
ecosystem as a unified network, discover latent structural
relationships, or support graph-based analyses of the system.

The Listening phase also presents a significant scalability challenge.
As the number of interviews, quotations, documents, and multimedia
sources increases, manually extracting perceptions, challenges,
opportunities, and other qualitative labels becomes increasingly
time-consuming and difficult to maintain consistently.

Furthermore, many of these labels rely heavily on human interpretation.
Different analysts may assign different perceptions or challenges to the
same piece of evidence, and the platform currently provides no
quantitative metrics to evaluate the confidence, consistency, or quality
of these annotations.

Another limitation is the incompleteness of the available data. Many
entities lack structured attributes such as geographic location,
temporal information, funding, budgets, or other contextual metadata,
reducing the ability to perform comprehensive quantitative analyses.

Most importantly, the platform lacks a mathematical representation of
the ecosystem. Relationships between actors, projects, narratives,
perceptions, and other conceptual layers are stored as database records
rather than mathematical objects. As a result, the platform cannot
directly leverage modern graph algorithms, network science, graph
machine learning, or large-scale semantic reasoning, all of which are
essential components of an accurate digital representation of this kind
of ecosystem.

The key contribution is not a single model but a pipeline that makes the
underlying data usable for graph analysis, narrative interpretation, and
controlled synthetic experimentation. The dataset contains initiatives,
agents, channels, perceptions, quotes, thematic objects, and relation
records. Those objects are not meaningful in isolation; they matter
because of how they connect.

The current 173 analysis shows a network with real structure but
substantial fragmentation. The source data contains roughly 350 records;
after graph construction, claim extraction, and normalisation, the final
graph contains 497 nodes and 654 edges, distributed across 204
connected components. The largest connected component contains 287
nodes (58 percent of the network). The system also contains 201
isolated nodes, 45 articulation points, and 61 bridges. These results
indicate that a relatively small number of nodes carry a
disproportionate share of the network's connectivity.

Perception diagnostics add a complementary interpretive layer. One
perception appears robust, while several others have low purity, low
support, or underdeveloped evidence. This suggests that the narrative
layer contains meaningful structure, but that the interpretive
categories are unevenly supported across the dataset.

The synthetic 173 layer is a controlled experimental extension. It keeps
the same topology as the original 173 dataset but adds invented budget
and investment values so that financial questions can be explored
without treating the synthetic values as empirical truth. In the present
report, this layer is used as a research sandbox rather than as evidence
of real funding conditions.

The dashboard integrates all outputs into a single Streamlit
application with tabs for the decision brief, overview map, geography
and evolution, network layers, health checks, listening data,
AI-generated semantic links, story clusters, narrative space,
perceptions, claims, GNN predictions, narrative simulation, structural
change, and (for synthetic platforms) budget and finance.

\section{Methods}

The entire pipeline is written in Python, with dashboard elements in
Streamlit and Plotly. Each subsection below describes one analytical
step, then shows the concrete results.

\subsection{Data recovery and normalization}

\textbf{Method.} KTool stores its data as nested records — a single project
entry contains embedded references to people, organisations, and
perceptions all mixed together. To analyse this data with graphs and
algorithms, we first need to untangle it into clean, separate tables.
The pipeline starts from the cleaned export of data hosted on the KTool
backoffice, and reorganises it into structured tables. This is necessary
because the source database stores relationships in nested,
hidden structures — a single project record may contain embedded
references to people, organisations, perceptions, and thematic tags, all
mixed together. If we flatten these relations too early, we would lose
the connections between them. The pipeline preserves the original
structure during recovery, then normalises the data into separate entity
tables: one for each kind of record in the ecosystem — agents, projects,
pilots, channels, quotes, perceptions, values, challenges, and thematic
areas. Normalisation makes it possible to join tables reliably, validate
completeness, and run the same analyses repeatedly. Because every table
carries a platform identifier, the pipeline can be run on any platform
at any time, making it straightforward to compare how a platform's
ecosystem evolves, or compare multiple platforms.

(Simple explanation: the pipeline takes messy, nested data from KTool
and organises it into clean separate tables — one for people, one for
projects, one for quotes, and so on — so that each type of information
can be analysed independently and then recombined.)

When the KTool source data has gaps such as missing locations, budgets,
descriptions, or dates, we can fill them with plausible values. For the
synthetic dataset, the pipeline assigns Irish town and county names
drawn from a curated list of 84 towns mapped to 28 counties, together
with realistic latitude/longitude coordinates, so that geographic
analyses can run even on records that arrived without location
information. Every filled value is tagged with a provenance label, so
the reader can always distinguish original data from synthetic data.

\textbf{Example.} For platform 173, the raw data contains 497 nodes across 12
types; after normalisation all 497 nodes have at least a label and a
node type, and 654 edges are recovered. For the synthetic dataset, the
pipeline assigns county and coordinate data so geographic analysis is
possible.

\textbf{Limitations.} The normalisation step depends on the completeness of
the source export. If relation tables are missing or sparsely populated,
the normalised output will undercount edges. The mapping of nested KTool
fields into flat tables is deterministic but not exhaustive: some
obscure field types in the source schema may be dropped silently.

\subsection{Graph construction}

\textbf{Method.} KTool's original format stores relationships as database
links between records — but an ecosystem is not a set of tables; it is a
network of actors, projects, and narratives that influence each other.
To capture this, we convert the normalised data into a graph where nodes
represent entities and edges represent relations between them.

The first family is declared relational links, which come directly from
the source database: an agent is linked to a project because they
participate in it, or a project is linked to a perception because it was
tagged with it. These 215 edges form the structural backbone for
platform 173.

(Simple explanation: these are the connections that users explicitly
entered into KTool — "this person works on this project." They are the
most trustworthy because they come directly from the source data.)

The second family is listening links, which trace the evidence chain
from a communication channel to an information item (a quote or
observation) to the values expressed in that item. These 136 edges
capture who said what and what values it reflected.

The third family is qualitative narrative links, which connect quotes
through semantic similarity, causal language, contradiction, and shared
themes. These 191 edges exist only for the synthetic dataset because
they require the quote extraction step that runs only on enriched data.

The fourth family is narrative claim edges, produced by the three-layer
claim extraction pipeline. Each claim node is connected to the quote
that generated it, and to any operational entities (agents, projects)
identified as subject or object of the claim. These 307 edges add a
structured semantic layer on top of the raw text.

The fifth family is interpretive edges, which are AI-inferred relations
between nodes that share thematic or semantic overlap. These 21 edges
are the most speculative and are clearly labelled so they can be excluded
from any analysis by selecting "Source only" in the dashboard.

This layered design avoids collapsing all evidence into one generic
network. It preserves the origin of each connection and makes it
possible to compare the ecosystem from different analytical angles. Each
edge family can be toggled on or off in the dashboard, and every view
can be further filtered by time window, location, value dimension, and
node type.

The sidebar provides dropdowns for time range (based on best available
date per node) and node type, a radio button to toggle AI-generated data
on or off, and a text input to switch platforms. Selecting "Source only"
removes any edge or node flagged as AI-generated (narrative claims and
qualitative narratives), leaving only declared relational, listening,
and interpretive edges.

\textbf{Example.} For the synthetic dataset, the full graph has 498 nodes and
870 edges across five families: 420 declared relational, 303 narrative
claim, 191 qualitative narrative, 136 listening, and 21 interpretive.
Toggling "Source only" removes 494 AI-generated edges and leaves 577,
revealing the structural core that exists without any AI inference.

(Simple explanation: of the 870 connections in the synthetic dataset,
494 were created by AI (claims and narrative links). Switching to
"Source only" hides these and shows only the 577 connections that
existed in the original data.)

With 870 edges among 498 nodes, the edge density is approximately 0.004,
meaning most node pairs have no recorded relationship regardless of
family. The interpretation of centrality and community metrics on such a
sparse graph should be treated as indicative rather than precise.

\textbf{Limitations.} The layered design preserves edge provenance but does
not address the fundamental sparsity of the graph. The interpretation of
centrality and community metrics on such a sparse graph should be treated
as indicative rather than precise.

\subsection{Structural diagnostics}

\textbf{Method.} A graph can look busy without being meaningful. Before adding
narrative layers, simulations, or predictions, we need to know whether
the network has real structure worth analysing — or whether it is just a
handful of isolated records. The pipeline runs a set of health checks on
the raw graph to answer basic questions about the ecosystem's shape.

The first check counts connected components. A component is a group of
nodes that are connected to each other through any path; nodes in
different components cannot reach each other at all.

(Simple explanation: imagine the graph as a set of islands. A connected
component is one island — you can walk from any node to any other node
on that island. If there are many small islands, the ecosystem is
fragmented.)

The second check measures edge density: how many connections exist
relative to what would be possible. If every node were connected to
every other node, density would be 1.0. Real ecosystems are usually
sparse, with density below 0.01.

The third check identifies orphan nodes that have no relationships at
all — zero connections to any other node. These are inactive records,
incomplete extractions, or genuinely disconnected entities.

The fourth check computes degree distributions and identifies the most
connected nodes using betweenness centrality, which measures how often
a node lies on the shortest path between other pairs of nodes. A node
with high betweenness acts as a gatekeeper — information flowing from
one part of the network to another is likely to pass through it.

(Simple explanation: betweenness centrality answers the question "if
this person disappeared, how many conversations would have to find a new
route?" A high score means the node is a critical bridge.)

The pipeline then runs a robustness simulation. It removes the top 30
nodes by betweenness centrality one at a time and measures how fast the
largest connected component shrinks, compared to removing 30 random
nodes. If targeted removal shrinks the network much faster than random
removal, the ecosystem depends heavily on a small set of hubs.

\textbf{Example.} Platform 173's graph contains 204 components. The largest
holds 287 nodes (58 percent of the network), meaning most activity is
concentrated in a single cluster while the remaining 210 nodes are
scattered across small isolated groups. The graph has a density of 0.005
(0.5 percent of all possible edges are present), which is typical for a
real-world ecosystem — sparse but structured. There are 201 isolated
nodes, giving an orphan rate of 40 percent. Rethink Ireland has the
highest betweenness centrality (0.12), followed by a handful of
Dublin-based organisations. Removing the top 30 nodes by betweenness
centrality shrinks the largest component by 34.2 percent, versus 12.9
percent for random removal — nearly three times as fast, confirming that
the network depends heavily on a small set of gatekeepers.

\textbf{Limitations.} Structural diagnostics describe the current state of
the graph but do not distinguish between missing data and genuine
disconnection. An orphan node may be an isolated actor or simply a
record that was never linked. The robustness simulation measures decay
under idealised conditions (sequential removal of the most central
nodes) and does not model real-world failure modes such as simultaneous
collapse or cascading defunding.

\subsection{Geographic analysis}

\textbf{Method.} From a governance perspective, knowing where activity happens
is only the first question. The more useful question is what kind of
activity happens where, and whether the geographic distribution of
attention matches the distribution of need or opportunity.

Geographic data comes from two sources. For the real 173 dataset, nodes
that have a location field are mapped directly. For the synthetic
dataset, a location enrichment script assigns Irish town and county
names together with latitude/longitude coordinates to projects, agents,
and channels that lack them. A separate enrichment step assigns each
node a county name by looking up its location in a hand-curated mapping
of 84 Irish towns to 28 counties, and backfills a value dimension from
the node's thematic areas using a rule-based keyword list of 7
dimensions and 34 keywords. The same script adds a simulated year column
for time-lapse animation, spreading nodes deterministically across years
2019 to 2025 using a hash of the node's global ID. This allows the
time-lapse view to show how value priorities shift over time even when
only one year of real dates is available.

(Simple explanation: if a record doesn't have a location, the pipeline
assigns one from a list of Irish towns so it can still appear on the
map. It also assigns a fake year so the time-lapse animation works.)

On the dashboard, the geography tab shows every mapped entity as a point
on an Open Street Map. The default view colours each point by entity
type (project, agent, channel, quote). Switching to the concentration
view scales each marker by the number of entities at that location,
revealing the urban bias of the ecosystem: Dublin accounts for roughly a
quarter of all mapped nodes, followed by Cork, Limerick, and Galway.
Rural counties such as Leitrim, Longford, and Carlow appear as isolated
points.

The underserve view goes a step further by calculating a simple support
ratio: how many agents and channels exist per project in each location.
A low ratio suggests a place where projects outnumber the organisational
infrastructure to support them.

(Simple explanation: if a county has many projects but few organisations
or channels to support them, it is "underserved" — there is activity but
not enough infrastructure.)

Two additional views connect geography to the value framework. The
county heatmap shades each Republic of Ireland county by its dominant
value dimension, using scattermap markers at county centroids with the
GeoJSON boundary file (26 counties from OpenStreetMap data). The
time-lapse view animates the same county-level data across the years
2019 to 2025.

\textbf{Example.} On the Geography tab, switching from the default point view
to the county heatmap colours Dublin dark blue (collaboration), Cork
green (social justice), and Galway orange (cultural identity), showing
how value orientations cluster spatially. Dublin accounts for roughly a
quarter of all mapped nodes. The underserve view highlights rural
counties where infrastructure lags behind project activity.

\textbf{Limitations.} The location-to-county mapping is a hand-curated list of
84 towns; any node with a town name outside this list is left
unassigned. The simulated year values are deterministic hashes, not real
dates, so the time-lapse shows structural change potential rather than
observed temporal trends. The county boundary GeoJSON (26 ROI counties)
excludes Northern Ireland counties entirely.

\subsection{Value dimensions and financial analysis}

\textbf{Method.} Two projects can both be about "youth" yet have completely
different motivations — one might be about fairness, the other about
innovation. To understand what is really driving an ecosystem, we need
to look beyond surface topics and map the deeper values behind them. A
thematic area such as "youth support" or "community inclusion" describes
what a project is about. A value dimension describes the deeper
rationale: whether the work is primarily about fairness
(social justice), about coordination (collaboration), about new
approaches (innovation drive), about evidence and measurement (evidence
based), about identity and heritage (cultural identity), about local
control (community autonomy), or about doing more with less (austerity
and scarcity). This distinction matters for decision making because two
projects with the same thematic tag can express very different value
orientations, and those orientations shape how stakeholders perceive
success.

(Simple explanation: the thematic area says what a project does — "youth
support." The value dimension says why it matters — is it about fairness,
innovation, tradition, or something else? Two projects doing the same
thing can have very different motivations.)

The pipeline assigns every project, quote, and perception one of the
seven value dimensions by mapping its thematic areas through a rule-
based lookup of 34 keywords. The mapping is transparent and can be
inspected in the enrichment script. In the dashboard, the Budget \&
Finance tab displays the result as a set of metric cards, one per value
dimension, showing the total investment in euros, the number of claims
and entities associated with that dimension, and the mean financial
bias.

For the synthetic dataset, the financial bridge reveals a clear
hierarchy: cultural identity receives 17.6 million euros, collaboration
10.1 million, innovation drive 7.0 million, social justice 5.2 million,
and evidence based 0.8 million, while community autonomy and austerity
and scarcity have zero recorded investment. For platform 173, no
financial information is available; the value dimension enrichment runs
but produces estimates rather than real budgets.

\textbf{Example.} The Budget tab's metric cards show cultural identity with
17.6 million euros and 20 claims, while community autonomy shows 0 euros
and 3 claims — a visible gap between narrative presence and financial
attention.

\textbf{Limitations.} The value dimension mapping is keyword-based and has no
disambiguation. A project tagged "economy" maps to innovation drive even
if the work is about poverty reduction. The financial data for the
synthetic dataset is pipeline-generated and has no empirical basis. For
the real 173 dataset, no budget information exists at all, so the
financial tab is disabled.

\subsection{Quote extraction and semantic linking}

\textbf{Method.} Most of what we know about an ecosystem comes from what
people say — interviews, focus groups, comments, observations. But raw
quotes are unstructured and hard to compare at scale. The listening layer
starts by extracting relevant quotes from the information table, scoring
them by quality, and detecting the semantic relationships between them. deliberately simple: length adds up to 0.35, action markers
add 0.25, value markers add 0.20, contradiction markers add 0.20, and
each structural link (channel, perception, theme) adds 0.10, all capped
at 1.0. Quotes scoring below 0.55 are excluded.

(Simple explanation: each quote gets a score from 0 to 1 based on how
long it is, whether it uses action words or value-laden language,
whether it contains contradictions, and whether it is connected to a
channel or theme. Only the most substantive quotes — those scoring above
0.55 — are kept.)

Once the quote candidates are selected, the pipeline detects semantic
relationships between them using TF-IDF cosine similarity combined with
linguistic markers. Every pair of quotes from different channels is
compared. If their cosine similarity exceeds 0.15 or they share a
thematic area, an edge is created. The edge type is determined by marker
detection: causal markers (because, therefore, enables, causes, etc.)
produce a causality edge, contradiction markers (but, however, despite,
barrier, etc.) produce a contradiction edge, and neither produces a
similarity edge. Edges with shared themes but low similarity (below 0.3)
are classified as frequency edges. Within the same channel, quotes are
linked in sequence edges ordered by their information ID, capturing the
flow of conversation.

(Simple explanation: quotes are compared to each other. If they use
similar language or share a theme, they get linked. If both use causal
words like "because" it is labelled a causality link; if one contradicts
the other, it is a contradiction link; otherwise it is a similarity
link. Quotes from the same channel are linked in order of conversation.)

\textbf{Example.} For the synthetic dataset, 72 quotes pass the threshold
with a mean score of 0.83. The pipeline creates 191 semantic edges: 57
sequence edges within channels, 54 causality edges, 40 contradiction
edges, 34 frequency edges, and 6 similarity edges. Every edge is tagged
with the edge family "qualitative\_narrative" and can be toggled on or
off in the dashboard. On the AI-Generated Links tab, selecting
"causality" filters the map to 54 purple edges connecting quotes like
"funding cuts led to service closures" and "short-term contracts create
instability."

\textbf{Limitations.} The marker-based edge detection captures only explicit
linguistic signals (because, therefore, but, however). Indirect causality
(a quote implies a cause without using causal language) is missed. The
TF-IDF similarity threshold of 0.15 is arbitrary; raising it would
increase precision but reduce recall. This step currently runs only for
platforms with enriched quote data (the synthetic dataset), because the
real 173 data lacks the text coverage needed for meaningful cross-channel
comparison.

\subsection{Three-layer claim extraction}

\textbf{Method.} The purpose of narrative extraction is to surface what people
in the ecosystem are saying and thinking, so that analysts can understand
the opinion landscape at a glance and use it to create or refine
perception profiles. Perceptions themselves are human-made constructs
that encapsulate population viewpoints; the pipeline's role is to make
the evidence behind them legible, quantifiable, and comparable.

(Simple explanation: we want to know what people are saying and what
values those statements reveal. This helps analysts build better
perception profiles — descriptions of what different groups in the
ecosystem believe.)

The extraction step processes every node of type information (quotes
from listening channels) that has at least 10 characters of text. For
each such node, a dependency parser converts the sentence into a
subject-verb-object triple. The parser used is hyperbase alphabeta, a
formal grammar engine, not a large language model: it applies a
predefined set of linguistic rules to identify the main predicate and
its arguments. This is deliberately conservative. It avoids the
hallucination risk of generative models and produces claims that are
directly traceable to the source text. The trade-off is coverage: quotes
with complex syntax, incomplete sentences, or non-standard grammar may
fail to parse and produce no claim at all.

(Simple explanation: the parser reads each quote and extracts the basic
structure: who did what to whom. It uses grammar rules, not AI, so the
results are predictable and traceable — but it will miss complex or
informal language.)

Alongside each surface claim, the pipeline checks for implicit meanings
using four linguistic markers. Negation markers (not, never, lack of,
fails to) invert the surface assertion: "we do not have funding" becomes
"system lacks funding." Emergency markers (crisis, at risk,
unsustainable) recast the statement as a need for protection: "youth
services are at risk" becomes "youth services need protection."
Conditional markers (if, unless, provided that) reframe it as a
dependency. Emotion markers capture urgency, frustration, hope, and fear
through a curated lexicon. If a quote triggers one or more of these
markers, an implicit claim is created alongside the surface claim. These
inference rules are rigid by design — negation always becomes "lacks",
emergency always becomes "needs\_protection" — which makes them
predictable and auditable but means they miss contextual nuance such as
sarcasm.

Every claim is then classified into one or more of the seven value
dimensions using a keyword-based lookup. Each dimension has 15-20
indicative terms. Every keyword match increments a score, and the
dimension with the highest score is assigned. The approach is
transparent but has no disambiguation: a sentence about "cutting-edge
research" matches both innovation drive and evidence based, and a
sentence about "budget cuts" matches austerity and scarcity even if the
speaker is criticising those cuts. The pipeline records the score and
assigns the highest-scoring dimension without applying a threshold.

Once claims are extracted, they are linked to the entities they mention.
Entity linking uses a composite scoring system. First, named entities
and noun phrases are extracted from the claim text using spaCy's
named-entity recognition model and noun chunker, plus a fallback
capitalised-phrase detector. Each mention is then scored against every
known operational entity in the graph (agents, projects, pilots, and
prototypes) using a weighted combination of exact label match (score
1.5), substring match (score 1.0), token overlap (Jaccard index scaled
to 1.2), embedding similarity from a sentence transformer model
(all-MiniLM-L6-v2, 384 dimensions, scaled to 0.5), and a graph context
bonus of 0.2 if the entity already has an edge to the source node. Links
with a composite score above the 0.6 threshold are kept. The threshold
was chosen heuristically and has not been empirically calibrated, which
means both false positives and false negatives are possible.

(Simple explanation: the pipeline checks each claim against known
organisations and people in the ecosystem. If a quote says "Dublin Youth
Collective needs funding," the pipeline tries to match "Dublin Youth
Collective" to a real node in the graph using name matching, word
overlap, and meaning similarity. Only confident matches are kept.)

The linked claims and entities produce three kinds of narrative claim
edges. First, a directed edge connects the source node (quote) to the
claim node it generated. Second and third, directed edges connect the
claim node to the subject entity and the object entity identified during
entity linking. In the synthetic dataset, these produce 303 edges, all
tagged with edge family "narrative\_claim" and with the provenance label
"is\_ai\_generated", so they can be toggled off in the dashboard.

Pattern analysis groups claims by their source node. If a single quote
generates claims mapped to more than one value dimension, it is flagged
as internally contradictory. For the synthetic dataset, several quotes
produce claims in both social justice and austerity and scarcity,
indicating contested discursive spaces. These patterns help analysts see
which perceptions may need splitting, merging, or refinement.

\textbf{Example.} For the synthetic dataset, the parser processes roughly 300
narrative nodes and produces 122 claims, which generate 303 narrative
claim edges. For a quote about "youth services funding," the hyperbase
parser extracts the surface claim "services support youth" (subject =
services, verb = support, object = youth). The negation marker on "lack
of funding" generates the implicit claim "system lacks funding." The
metanarrative classifier assigns it to both social justice (keywords:
youth, inclusion) and austerity and scarcity (keyword: funding), flagging
it as contradictory on the dashboard's Claims tab.

\textbf{Limitations.} The hyperbase parser has limited coverage of informal
speech, incomplete sentences, and code-switched language. The entity
linking threshold of 0.6 was chosen heuristically, so both false
positives and false negatives are possible. The implicit claim detection
captures only the four marker types and misses rhetorical questions,
analogy, or understatement. All AI-generated edges are labelled but
downstream metrics treat them identically to source edges unless the
user selects "Source only."

\subsection{Story clustering and narrative profiles}

\textbf{Method.} Individual claims are too numerous to inspect one by one.
Grouping them into related clusters reveals the main conversation threads
running through the ecosystem — which topics are prominent, which are
marginal, and how different stakeholder groups talk about the same issue
from different angles. The quote clusters produced in this step are the
closest the pipeline comes to a "narrative space" — a map of what people
in the ecosystem are talking about and how those conversations group
together. These narrative spaces are the raw material analysts use to
construct perception profiles: each cluster represents a candidate
viewpoint that may warrant its own perception, and the claims within it
reveal the value dimensions driving that viewpoint.

(Simple explanation: think of each cluster as a conversation group.
Everyone in the same cluster is talking about similar things. Analysts
can look at a cluster and ask: "does this feel like one coherent
viewpoint, or are there multiple opinions mixed together that should be
separated into different perception profiles?")

After semantic edges are created, the pipeline clusters quotes into
story groups using TF-IDF vectorisation and agglomerative clustering.
The number of clusters is set dynamically based on quote count, ranging
from 2 to 8. For the synthetic dataset, 72 quotes are grouped into 8
clusters.

Each cluster receives a label derived from its most frequent thematic
area, a representative quote (the one with the highest selection score),
and a profile of the channels and themes it covers. Across the 8
clusters, the average selection score ranges from 0.68 to 0.95.

These clusters are then linked to the claims extracted in the previous
step by matching the claim's source node ID to the information ID in the
cluster table. The resulting cluster-claim matrix shows each cluster's
value dimension profile. If a single cluster contains claims in more
than one value dimension, it is flagged as internally contradictory. For
the synthetic dataset, some clusters span both social justice and
innovation drive claims. This flags for the analyst whether a cluster
represents a single coherent perception or an amalgamation that may need
splitting.

A separate narrative diffusion analysis computes PageRank for every node
in the full graph and a diffusion bias score that contrasts personalised
PageRank between the two largest agent types. The top-ranked node by
PageRank is typically the most connected agent in the mapping layer,
while the bias score reveals which voices carry further in one agent
network versus another.

(Simple explanation: PageRank measures which voices are most influential
in the network. The bias score shows whether certain types of
organisations have more reach than others.)

\textbf{Example.} Cluster 3 ("youth engagement") contains 14 quotes about
education and social inclusion, with claims split between social justice
(8 claims) and innovation drive (3 claims). It is flagged as internally
contradictory. On the dashboard's Story Clusters tab, the cluster's
representative quote is shown alongside its value dimension profile.

\textbf{Limitations.} Clustering uses TF-IDF similarity, which captures word
overlap but not synonymy or discourse framing. Clusters with fewer than
three quotes have unreliable profiles. The number of clusters (between 2
and 8) is set by a heuristic formula, not validated against a gold
standard. The narrative diffusion PageRank is computed on the full graph
including AI-generated edges unless filtered.

\subsection{Opinion simulation (Friedkin-Johnsen)}

\textbf{Method.} Knowing what people say in isolation tells us the topics, but
ecosystem analysis is also about influence — how opinions spread from
one group to another, which viewpoints are likely to reinforce each
other, and which are structurally isolated from the rest of the
conversation. Simulating these dynamics on the cluster network reveals
influence pathways that are invisible when clusters are examined
independently. The opinion simulation treats each story cluster as a
node in a smaller network. A cluster's opinion is a seven-dimensional
vector, one entry per value dimension, built by counting how many of its
claims fall into each dimension and normalising the result. Two clusters
are connected if a semantic similarity edge exists between any quote in
the first cluster and any quote in the second. In the synthetic dataset,
several clusters are connected through shared thematic areas even though
they express different value dimensions, creating a cross-cluster
influence network.

(Simple explanation: each cluster has an "opinion profile" — for example,
70 percent social justice, 30 percent innovation drive. Clusters that
share themes can influence each other's opinions.)

The simulation uses the Friedkin-Johnsen model, a well-established
framework for how opinions evolve in a connected group. Each cluster has
an anchoring opinion (its own initial profile) and is influenced by the
opinions of the clusters it is connected to. The balance between
anchoring and influence is controlled by a stubbornness coefficient set
proportional to the cluster's claim count: clusters with 15 or more
claims are highly stubborn (stubbornness near 0.9), while clusters with
only 2-3 claims are highly persuadable (stubbornness near 0.2). The
model iterates until the opinions stabilise, typically within 20-50
iterations.

(Simple explanation: clusters with lots of evidence (many claims) are
stubborn and hard to sway. Clusters with little evidence are easily
influenced by their neighbours. The model runs until everyone settles
into a stable opinion.)

The baseline answers the question: if these story clusters were left to
interact naturally, what would each cluster's final opinion profile look
like? The dashboard visualises this as a radar chart with seven axes,
one per value dimension.

From the baseline, the pipeline runs intervention simulations. For each
cluster in turn, it adds an external agent with a neutral opinion (equal
weight across all seven dimensions) and a low stubbornness of 0.2,
connects it to that cluster, and re-runs the simulation. The question
is: if someone new entered the ecosystem and engaged primarily with this
one cluster, how would every other cluster's opinion shift? For the
synthetic dataset, the largest opinion shifts occur when connecting the
agent to the most central cluster (by cross-cluster edge count), with
max deltas of approximately 0.02 to 0.05 in the affected dimensions.

(Simple explanation: the simulation answers "if a new person started
talking to group A, how would groups B, C, and D change their minds?"
It is a way to test which conversations have the most ripple effects.)

A second, more detailed simulation (GBCM-FJ) incorporates perception
nodes directly, using their embedding similarity, value dimension
correlation, and claim counts to compute influence. This model runs on
the full graph (not just clusters) and evaluates counterfactual edges
from the GNN link prediction output. It computes how disagreement,
diversity, and perception PageRank change when a proposed link is added.

\textbf{Example.} The Narrative Simulation tab shows an 8-cluster radar chart.
Cluster 1 ("community development") has high social justice and low
innovation drive. Connecting a neutral agent to it shifts Cluster 4
("digital inclusion") toward social justice by 0.04 on the normalised
scale — a small but detectable influence cascade.

\textbf{Limitations.} The stubbornness formula is proportional to claim count,
making clusters with few claims highly persuadable regardless of whether
those claims are strongly held. The model has no temporal dimension. The
intervention tests only a neutral agent, not a funder, advocate, or
traditionalist. Semantic edges between clusters capture textual
similarity, not organisational ties.

\subsection{Perception diagnostics}

\textbf{Method.} Perceptions in KTool are human-made profiles that encapsulate
the opinions of a population segment. The pipeline does not generate them
automatically. Instead, it evaluates each existing perception on five
health metrics so analysts know which are well-supported by evidence and
which need revision.

(Simple explanation: perceptions are descriptions of what different
groups in the ecosystem believe — for example, "young people in Dublin
prioritise job opportunities." The pipeline checks each perception
against the actual quotes to see if the evidence supports it.)

Internal coherence measures the mean cosine similarity of semantic edges
between quotes within the same perception. A high coherence means the
quotes are all saying similar things. Purity measures what fraction of
each quote's top-5 semantic neighbours belong to the same perception. A
high purity means the perception's quotes form a clean cluster. Source
entropy measures the diversity of channels feeding each perception —
a perception fed by many different channels is more robust than one
from a single source. Contradiction density measures the fraction of
intra-perception edges that are contradictions — a high value means the
perception contains conflicting views. Each perception also receives a
status flag: "Robust," "Underdeveloped (< 3 quotes)," "Low coherence,"
"Single channel," "High internal contradiction," or "Low purity."

(Simple explanation: the five metrics answer: are the quotes consistent?
Do they belong together? Do they come from diverse sources? Do they
contradict each other? A "Robust" perception passes all checks.)

The claim extraction output feeds directly into this diagnostic layer:
an analyst reviewing a perception with "Low purity" can drill into the
claims behind its quotes to see which value dimensions are pulling it in
different directions, then decide whether to split it into two more
coherent profiles or merge it with a neighbouring perception.

A second phase evaluates narrative alignment: how distinct each
perception's value dimension profile is from the others. The pipeline
computes a silhouette score (how far a perception's claim profile is
from its nearest neighbour) and flags narratively redundant pairs (claim
cosine similarity above 0.80). This helps identify perception categories
that may be analytically indistinguishable and candidates for merging.

\textbf{Example.} For platform 173, the diagnostics reveal one robust
perception with high coherence and moderate quote count, alongside
several perceptions with fewer than 3 quotes or low internal coherence.
On the Perceptions tab, "community wellbeing" has an internal coherence
of 0.72 (high), a purity of 0.55 (moderate), and a silhouette score of
0.31 — distinct from "youth development" but not sharply separated.

\textbf{Limitations.} Perception diagnostics depend on the availability of
semantic edges between quotes. If the quote set is small or the semantic
edge detector produced few edges, the coherence and purity estimates are
unreliable. The status flags use heuristic thresholds (coherence below
0.40 is "low," purity below 0.30 is "low") that have not been validated
against human judgement.

\subsection{Link prediction and project recommendations}

\textbf{Method.} The graph contains only the relations recovered from source
data, but a government or funding agency can actively create new
connections: fund a project in a new region, broker a partnership between
two organisations, or convene stakeholders around a shared challenge. The
pipeline uses a graph neural network to identify which of these possible
new connections would have the largest structural and narrative impact,
and then ranks them by practical feasibility for decision-makers.

(Simple explanation: the graph only shows connections that already exist.
But funders and policymakers can create new connections — for example,
introducing a youth organisation to a potential funder. The GNN predicts
which new connections would be most valuable.)

The first step is preparing the graph for the model. Every node is given
a feature vector that combines its position in the network (degree,
centrality, k-core level, whether it is an articulation point), its type
and phase (one-hot encoded), its textual description (converted to a
fixed-size hash vector of 256 dimensions), and its semantic embedding
from the sentence transformer model (384 dimensions). Every edge is
given an omega confidence score: omega = 0.45 x provenance (how the edge
was created) + 0.35 x alignment (cosine similarity of endpoint
embeddings) + 0.20 x variance stability (how consistent the alignment
is). Edges with omega below 0.6 are excluded from training. The GNN
directory also records feature slices (numeric, categorical, text hash,
semantic) so the contribution of each feature type is traceable.

The link prediction model is a graph autoencoder. A two-layer GCN
encoder compresses each node into a 64-dimensional embedding. A dot-
product decoder takes any pair of node embeddings, computes their dot
product, and passes it through a sigmoid function to produce a
probability between 0 and 1. The model is trained on 70 percent of the
existing edges, with 15 percent for validation and 15 percent for
testing.

(Simple explanation: the model learns to compress each node into a
compact representation (64 numbers) that captures its role in the
network. Then for any two nodes, it asks: "given their roles, should
they be connected?" It outputs a probability from 0 to 1.)

Once trained, the model scores all possible pairs of nodes that are not
already connected. Only pairs where both endpoints belong to the mapping
layer agents and projects are retained, because these are the entity
types where a new connection translates into an actionable
recommendation. The top 50 candidates are then scored by a structural
impact engine that answers: what would actually happen if this link were
added?

The impact engine temporarily adds each candidate edge to a copy of the
full graph and measures what changes. It runs a breadth-first search
from both endpoints before and after the addition, counting how many new
perception nodes and claim nodes become reachable within three hops. It
checks whether the link unlocks value dimensions that were previously
inaccessible. It classifies the bridge type: a learning bridge unlocks
entirely new value dimensions, a cleavage breach breaks through a
structural divide, a coalition reinforcement connects nodes that already
share values, and a structural link adds no new narrative pathways. It
also checks whether the link merges two separate connected components,
whether either endpoint is an articulation point, and how semantically
compatible the two nodes are.

(Simple explanation: adding a new connection can unlock new ideas and
perspectives by creating pathways to parts of the network that were
previously unreachable. The impact engine measures exactly what becomes
reachable and how valuable it is.)

All of these factors are combined into a composite governance score.
Narrative impact (new pathways and new value dimensions) carries the
most weight at 35 percent, followed by bridge type at 20 percent,
component merger at 15 percent, and degree, articulation-point status,
and domain compatibility at 10 percent each. Each candidate also
receives a feasibility badge: fundable (strong semantic alignment),
cross-domain (moderate alignment with structural value), or topological
(weak alignment, purely structural connection). The dashboard's GNN
Predictions tab displays the top recommendations sorted by this
composite score, with the badge, bridge type, number of new pathways,
and a rationale.

A separate scoring process identifies prototype project candidates from
the full set of project nodes. Each project is scored on a weighted
combination of its budget (32 percent), its betweenness centrality (28
percent), the number of perceptions it is connected to (16 percent), the
number of agents involved (12 percent), and the number of thematic
topics it covers (12 percent). The dashboard's Suggested Projects
section displays the top candidates with their estimated budget, score,
and a suggestion reason.

\textbf{Example.} On the synthetic dataset with 498 nodes and 1071 edges (1625
after bidirectional expansion and omega filtering), the model achieves a
best validation AUC of 0.986, but the test AUC drops to 0.72, confirming
overfitting on a small graph. The GNN Predictions tab shows a recommended
link between "Dublin Youth Collective" (agent) and "Creative Communities
Programme" (project) with a governance score of 0.81. Adding it would
unlock 3 new perception paths and bridge the social justice and cultural
identity value dimensions. The badge is "Fundable."

\textbf{Limitations.} The GNN achieves high validation AUC but on a small
graph, so the ranking is more reliable than the absolute probability
values. The narrative impact scores depend on a BFS horizon of three
hops, which may miss longer-range effects. The feasibility badge reflects
semantic compatibility, not organisational or political readiness. The
budget data for the synthetic dataset is pipeline-generated.

\subsection{Structural change feasibility}

\textbf{Method.} The question that motivates this section is not "what should
change?" but "given the current relational structure, what kinds of
change are even possible?" This is a fundamentally different question
from the link prediction in the previous section, and it is grounded in
political science rather than graph theory alone. The Advocacy Coalition
Framework (Sabatier and Jenkins-Smith, 1993) holds that policy
subsystems are organised around competing coalitions bound by shared
beliefs, and that major change requires external shocks or the
accumulation of minor perturbations over long periods. Punctuated
Equilibrium Theory (Baumgartner and Jones, 1993) describes how long
periods of policy stability maintained by dense, insular policy
monopolies are interrupted by brief bursts of rapid reconfiguration when
excluded actors reframe the policy image and find new institutional
venues. Path dependency theory (Pierson, 2000) warns that early
institutional choices and funding patterns lock ecosystems into
trajectories that become increasingly costly to reverse. These
frameworks provide the interpretive lens for the metrics below.

(Simple explanation: we use political science theories to interpret the
network structure. For example, a tightly connected core of organisations
might be a "policy monopoly" that resists change. Isolated nodes might be
"excluded voices" that could trigger change if they found new allies.)

The pipeline computes four composite dimensions from the network
metrics, each mapped to a political science construct.

The first is leverage. In network terms, nodes with high betweenness
centrality occupy structural holes and act as gatekeepers of information
flow. In political terms, they are policy brokers and boundary spanners
who can mediate across coalitions, build consensus, and accelerate
narrative diffusion (Gould and Fernandez, 1989). The pipeline identifies
the top betweenness nodes and top bridge agents as potential leverage
points. A high leverage score means the ecosystem has clear hubs that
can be strengthened to accelerate diffusion, but losing them would
fragment coordination.

(Simple explanation: leverage measures whether there are key
organisations that, if strengthened, could spread ideas more effectively.
Like identifying the best-connected people in a room — they already have
everyone's ear.)

The second is plasticity, which measures the network's spare capacity to
form new connections without collapsing existing structures. Nodes with
degree one or two are peripheral and easy to rewire; nodes with degree
three to five are candidates for new brokering roles. The count of
high-confidence GNN link predictions also contributes, because each
predicted link represents a latent connection that could be activated.
A high plasticity score means the periphery has abundant unused
connection potential — there are many isolated or weakly connected nodes
waiting to be integrated.

(Simple explanation: plasticity measures how much room there is to create
new connections. A high score means there are many isolated nodes that
could be brought into the network — like empty chairs at a table.)

The third is blockage. Fragile connectors are articulation points that
also carry bridge edges: removing them would both disconnect the network
and break a unique link. These are institutional bottlenecks, single
points of systemic failure. Isolated nodes have degree zero and
participate in no relationships at all. Small components with three or
fewer nodes cannot sustain meaningful diffusion. Blocked perceptions are
perceptions that have no path to any information node in the graph,
trapping them in discursive silos — a structural analogue of the
discursive cleavages and ideological polarisation described in
polarisation research.

(Simple explanation: blockage measures single points of failure — nodes
that, if removed, would break the network. Like a bridge that, if
destroyed, would cut off an entire island.)

The fourth is lock-in or path dependency. A network with a very dense
core (a high fraction of nodes in the maximum k-core layer) resembles a
policy monopoly: a dominant coalition controls the most connected
positions, and new entrants or alternative viewpoints struggle to gain
traction (Baumgartner and Jones's concept of the "policy monopoly"
maintained by a supporting policy image and an insulated venue). The
robustness gap — how much faster the network decays under targeted
attack versus random failure — reinforces this picture: if removing the
most central nodes causes disproportionate damage, the network depends
heavily on those hubs and is structurally brittle.

(Simple explanation: lock-in measures how stuck the network is. If a
small group of organisations controls all the connections, it is hard
for new voices to break in. Like a room where the same few people do
all the talking.)

These four dimensions are normalised to scores between zero and one and
combined into an overall readiness score, where blockage carries the
highest weight. The dashboard's Decision Brief tab displays all four
scores as horizontal bars with colour coding (green for favourable, amber
for moderate, red for concerning), each accompanied by a plain-language
explanation: leverage becomes "fraction of nodes that are well-connected
hubs," blockage becomes "fraction of nodes that are single points of
failure," and lock-in becomes "how stuck the network is."

From these scores, the pipeline generates named hypotheses grounded in
the political science frameworks above. A hypothesis about a high-
betweenness node cites policy brokerage and boundary spanning theory
(Gould and Fernandez). A hypothesis about the innermost k-core layer
cites the Advocacy Coalition Framework and the concept of a policy
monopoly (Sabatier and Jenkins-Smith; Baumgartner and Jones). A
hypothesis about blocked perceptions cites discursive cleavages and
ideological polarisation. Each hypothesis is assigned a confidence level
(high or medium) based on thresholds derived from the data distribution,
and only high-confidence hypotheses are promoted to recommendations. The
dashboard's Structural Change tab lists these hypotheses alongside their
framework reference, confidence level, and the specific node identifiers
they reference.

\textbf{Example.} For the synthetic dataset: leverage is 0.78 (moderate —
betweenness concentrated in Dublin organisations), plasticity is 1.0
(maximum — sparse periphery with ample rewiring capacity), blockage is
0.56 (moderate — some fragile connectors and small components), lock-in
is 0.16 (low — only 2 percent of nodes in the innermost k-core of 8).
Overall readiness is 0.74. The Structural Change tab shows leverage at
0.78 (green), plasticity at 1.0 (green), blockage at 0.56 (amber), and
lock-in at 0.16 (green). The caption reads: "Network is open to
reconfiguration. Blockages (fragile connectors, small components) need
attention."

\textbf{Limitations.} The change readiness scores are computed from a static
graph snapshot. Without temporal data, the pipeline cannot distinguish
between a network that is slowly rewiring and one that is frozen in
place. The political science mappings (betweenness = broker, dense core
= policy monopoly) are analogical, not proven by the data — they provide
an interpretive frame, not a causal test. The Gould-Fernandez brokerage
typology is not fully implemented (the pipeline uses simple betweenness
centrality, not group-attribute brokerage classification). The blockage
score weights each type of blockage equally, which may not reflect their
relative policy significance.

\subsection{Dashboard implementation}

\textbf{Method.} The dashboard is the main user-facing interface for the
outputs. It is built with Streamlit and Plotly and presents the analysis
through 14 or 15 tabs, depending on whether the platform has synthetic
financial data. The design goal is to make the analytical layers
understandable to readers who do not need to inspect code or raw tables.

(Simple explanation: the dashboard is a web app that shows all the
analysis results in one place, organised into tabs. No coding required.)

The sidebar provides platform selection, AI-data toggle (All data /
Source only), time range slider, and edge family multiselect. The main
area shows the network map, health checks, diagnostic tables, simulation
results, and recommendation lists.

\textbf{Example.} The sidebar's time filter defaults to the full date range
(2019 to 2025 for the synthetic dataset). Selecting "Source only" and
narrowing to 2023-2024 immediately reduces the graph from 498 nodes to
the subset with dates in that window, and removes all AI-generated edges
from the view.

\textbf{Limitations.} The dashboard reads pre-computed CSV and JSON files from
the analysis directory. If a script has not been run for the selected
platform, the corresponding tab shows a "data not available" message
rather than an error. Performance degrades for very large graphs (above
1,000 nodes) because the network visualisation uses a force-directed
layout computed in real time. The dashboard is designed for exploratory
analysis, not for reproducible reporting: it does not export figures or
log user interactions.

\section{Discussion}

The main substantive finding is that the ecosystem has meaningful
structure, but that structure is unevenly distributed. There is a clear
core, there are many bridges, and there are many isolated nodes. This
combination means the network is neither trivial nor fully integrated.
It is a functioning socio-semantic system with visible structural
dependencies.

From a research perspective, this matters in three ways. First, it shows
that graph analysis is appropriate for the dataset because there is
enough structure to study. Second, it shows that narrative and
perception layers add value because they reveal interpretive patterns
that are not visible in the graph alone. Third, it shows that the
current system still has data limitations, especially where relation
coverage is incomplete or where categories remain weakly supported.

The results also imply that interventions should probably focus on
improving relational completeness, strengthening bridge entities, and
standardizing how narrative evidence is linked to structured records. In
a practical sense, a better-connected data model would make later
analysis more reliable and more interpretable.

\section{Limitations}

The current report and pipeline are limited by the available source
data. Some relation families are only partially populated, and some
interpretive categories have weaker support than others. That means the
results should be read as a careful description of the present dataset,
not as a full account of the entire underlying ecosystem.

A second limitation is that the synthetic financial layer is
experimental. It is useful for scenario analysis, but it should not be
mistaken for real-world financial evidence. Any paper using this layer
should state that the budgets and investment values were invented to
support controlled testing.

A third limitation is that the current outputs are strongest for
description and structure. They are less able to support causal claims
about change, because the project is not yet based on a longitudinal
sequence of repeated snapshots. For that reason, the present report
should be treated as a foundation for future work rather than as a final
causal explanation.

\section{Conclusion}

The ALC K-Tool project provides a structured way to move from a raw
ecosystem record dump to an interpretable analytical system. The value
of the work is in the combination of recovery, normalization, graph
construction, narrative analysis, and dashboard presentation.

The Platform 173 results show a sparse and fragmented network with a
small but meaningful core, many isolated nodes, and substantial
sensitivity to targeted node removal. The perception results show that
the narrative layer contains both robust and weakly supported
categories. The synthetic 173 layer adds a controlled way to test
financial-style questions without changing the original topology.

Taken together, these results support a straightforward conclusion: the
dataset is now rich enough for serious descriptive and interpretive
analysis, and the pipeline is strong enough to support future research
on structure, narrative, and investment alignment. The most important
next step is not a more complicated model for its own sake, but better
data coverage, clearer relational labeling, and repeated observations
over time.

\section{Open questions and future research}

The pipeline and results presented in this report raise several questions
that the current data cannot fully answer. These are outlined below as
directions for future work.

\textbf{Static snapshot versus evolving system.} The entire analysis is based
on a single export of the KTool database. A graph built from one point
in time shows structure but cannot show how that structure formed,
whether it is stable, or which edges are new versus old. The most
important next step is to collect data at regular intervals — quarterly
or annually — so the pipeline can track changes as they happen. Without
repeated observations, the structural change feasibility section remains
a theoretical estimate, not a measured trend.

\textbf{Missing data and orphan nodes.} Across both datasets, roughly 40
percent of nodes have no connections. Some of these may be genuinely
isolated actors, but many are almost certainly records that were never
linked in the source database — for example, a project and its lead
organisation exist in separate tables but were never associated in
KTool's interface. The pipeline cannot distinguish between these two
cases. Future work should audit a sample of orphan nodes manually to
estimate the share of false orphans, and develop heuristics that recover
likely missing links before the graph is built.

\textbf{Validation of narrative extraction.} The claim extraction step
surfaces structured claims from quotes automatically. But there is no
ground-truth dataset that tells us whether the parser correctly
identifies the intent of a quote, or whether it generates false positives
(claims the speaker never intended). A validation study — where human
coders label a random sample of quotes and their extracted claims are
compared — would measure precision and recall. Without it, the narrative
statistics should be treated as exploratory, not confirmatory.

\textbf{Perception-to-claim alignment.} The pipeline scores perceptions on how
well they match the claim data, but it does not yet tell analysts how to
fix a weak perception — for example, whether to broaden its definition,
add supporting quotes, or retire it entirely. A future version could
generate specific revision suggestions: "Perception X scores low on
coverage because its definition mentions value dimension Y, but the
closest claim cluster expresses dimension Z. Consider rewording the
perception to match the available evidence."

\textbf{Simulation accuracy.} The Friedkin-Johnsen opinion simulation produces
convergence patterns and influence rankings, but these have not been
compared against real opinion change over time. A natural next step is
to collect survey data or interview transcripts from the same ecosystem
at two time points, run the simulation on the first time point, and
check whether the predicted trajectory matches the observed second time
point. Until then, the simulation is best understood as a what-if tool
rather than a forecast.

\textbf{Cross-platform generalisation.} The pipeline was developed and tested
on Platform 173 (real data) and Platform 173\_synthetic (synthetic
financial overlay). Other KTool platforms may have different field
structures, different data quality, or different entity types. The
recovery and normalisation stages are designed to be generic, but the
narrative extraction rules and value dimension taxonomy may need
adjustment for ecosystems in different sectors or languages. Replicating
the pipeline on two or three additional platforms would reveal which
parts generalise and which are platform-specific.

\textbf{Causality and policy impact.} The current analysis is descriptive and
diagnostic — it tells you what the ecosystem looks like, where it is
fragile, and which narratives exist. It does not tell you whether a
specific intervention (funding a new project, brokering a partnership)
causes a measurable improvement. Moving from description to causal
inference would require an experimental design: for example, using the
link prediction recommendations in a real funding cycle, then measuring
whether the recommended connections actually produced better outcomes
than non-recommended ones. This is a long-term goal, but the pipeline
already provides the measurement infrastructure needed to attempt it.

\section{Resources}

The following works are referenced in the methods and interpretation
above.

Baumgartner, F. R., and Jones, B. D. (1993). *Agendas and Instability in
American Politics*. University of Chicago Press.

Friedkin, N. E., and Johnsen, E. C. (1999). "Social influence networks
and opinion change." \textit{Advances in Group Processes}, 16, 1--29.

Gould, R. V., and Fernandez, R. M. (1989). "Structures of mediation: A
formal approach to brokerage in transaction networks." *Sociological
Methodology*, 19, 89--126.

Pierson, P. (2000). "Increasing returns, path dependence, and the study
of politics." \textit{American Political Science Review}, 94(2), 251--267.

Sabatier, P. A., and Jenkins-Smith, H. C. (1993). *Policy Change and
Learning: An Advocacy Coalition Approach*. Westview Press.

\end{document}